\section*{Chapter \ref{ch:b3:gravity} --- Gravity and the Trade Data}

\begin{solution}{exo:b3:gravity:1}
By $2^{1.18} - 1 = 127~\%$: the slope is 1.18. % figdata: opening_stats
At the left end, the binned points lie above the line: the smallest and most
distant pairs trade more than the line predicts, partly because only pairs with
positive trade are in the data --- among small, distant pairs, those that trade at
all are the ones with unusually large flows.
\end{solution}

\begin{solution}{exo:b3:gravity:2}
$X_A/X_B = (500/1{,}000) \times (1{,}000/4{,}000)^{-0.9} = 0.5 \times 4^{0.9} = 0.5 \times 3.482 = 1.74$:
pair A trades 1.74 times as much.
\end{solution}

\begin{solution}{exo:b3:gravity:3}
$2^{-0.9} - 1 = -46.4~\%$; $0.9^{-0.9} - 1 = +9.9~\%$.
\end{solution}

\begin{solution}{exo:b3:gravity:4}
Two large, close countries with no other obstacle between them that trade far
less than gravity predicts --- or far more than a pair of the same size and
distance. One pair can be explained by its own pair variables (an embargo, a war,
a shared language the regression omits); gravity is a statement about the
average across pairs, and is rejected only if the average relation fails.
\end{solution}

\begin{solution}{exo:b3:gravity:5}
The tested solution is \code{ex05_language.py}, in the chapter's \code{exercises/}
folder. Distance moves from $-1.206$ to $-1.175$ and colonial link from 1.412 to 0.945;
the language coefficient is 0.812, so a common language multiplies trade by 2.25.
% output: ex05_language
Former colonies often share their former ruler's language, so without the
language dummy the colonial link picked up part of the language effect; close
countries also share languages more often, so distance had absorbed a little of
it too.
\end{solution}

\begin{solution}{exo:b3:gravity:6}
The tested solution is \code{ex06_fta.py}, in the chapter's \code{exercises/}
folder. 2005: OLS 0.503, trade $\times 1.65$ (+65~\%); \omterm{def:b3:gravity:ppml}{PPML} 0.373, $\times 1.45$ (+45~\%). 2019:
OLS 0.464, $\times 1.59$; \omterm{def:b3:gravity:ppml}{PPML} 0.380, $\times 1.46$. % output: ex06_fta
\omterm{def:b3:gravity:ppml}{PPML} gives the smaller effect in both years; the gap narrows from 0.13 to 0.08,
but given standard errors of 0.04 to 0.06 the difference between the two years is
not clearly significant.
\end{solution}

\begin{solution}{exo:b3:gravity:7}
Not sharing a border multiplies trade by $e^{-b}$, equivalent to a tariff $e^{b/\varepsilon} - 1$:
$e^{1.027/5} - 1 = 22.8~\%$ with fixed effects by OLS, $e^{0.537/5} - 1 = 11.3~\%$ by \omterm{def:b3:gravity:ppml}{PPML}.
\end{solution}

\begin{solution}{exo:b3:gravity:8}
Trade varies as $t^{-5}$: $(1.2/1.1)^5 = 1.545$, a rise of 54.5~\%. With \omterm{def:b3:gravity:multilateral-resistance}{multilateral
resistance} the effect on this pair is smaller, since the country's resistance
falls as one of its costs falls; trade with other partners falls, because the
partner with the new port is now relatively cheaper (diversion).
\end{solution}

\begin{solution}{exo:b3:gravity:9}
The log of zero is not defined, so the rows drop out. The pairs that survive at
long distances are those with unusually large shocks; the regression sees
large flows where the truth has many zeros, so distance seems to cost less than
it does: the elasticity is biased toward zero.
\end{solution}

\begin{solution}{exo:b3:gravity:10}
The tested solution is \code{ex10_distance_puzzle.py}, in the chapter's
\code{exercises/} folder. The elasticities are $-1.345$, $-1.425$, $-1.453$, $-1.503$,
$-1.543$, $-1.506$ and $-1.440$; the trend is $-0.048$ per decade (s.e.\ 0.026), $t = -1.82$.
% output: ex10_distance_puzzle
If anything, distance mattered slightly more over the period, though the trend is
not significant at 5~\% (and seven estimates from overlapping data are a thin basis
for a test). The elasticity is not falling: the puzzle.
\end{solution}

\begin{solution}{exo:b3:gravity:11}
By symmetry $\Pi_1 = \Pi_2 = P_1 = P_2 = P$, and $P^{1-\sigma} = \tfrac12(1 + t^{1-\sigma})P^{\sigma - 1}$, so
$P^{2(1-\sigma)} = \tfrac12(1 + t^{1-\sigma})$. Then $X_{12}/X_{11} = t^{1-\sigma}$: imports relative to home
purchases are $t^{1-\sigma}$. The import share of spending is $X_{21}/(X_{11} + X_{21}) = t^{1-\sigma}/(1
+ t^{1-\sigma})$; it falls with $t$ and falls faster the larger $\sigma$, from $\tfrac12$ at $t = 1$
toward 0.
\end{solution}

\begin{solution}{exo:b3:gravity:12}
The tested solution is \code{ex12_ppml_vs_ols.py}, in the chapter's
\code{exercises/} folder. Mean OLS slope 0.490, mean \omterm{def:b3:gravity:ppml}{PPML} slope 0.996; spreads 0.080
and 0.128. % output: ex12_ppml_vs_ols
$E[\ln \eta \mid x] = -(0.5 + x)/2$, so $E[\ln y \mid x] = 0.75 + 0.5x$: OLS estimates
$1 - 0.5 = 0.5$, exactly what it finds. \omterm{def:b3:gravity:ppml}{PPML} is consistent, at the cost of a larger
spread.
\end{solution}

\begin{solution}{pb:b3:gravity:union}
\textbf{1.} OLS 374, \omterm{def:b3:gravity:ppml}{PPML} 380. OLS on logs drops the 6 pairs with no recorded
trade. % figdata: union_estimates, union_stats

\textbf{2.} $2^{-0.904} - 1 = -46.6~\%$.

\textbf{3.} $(300/3{,}544)^{0.680} \times (1.08/4.31)^{-0.904} = 0.187 \times 3.48 = 0.65$: Brenland should
export to Kestria about two thirds of what it exports to N10. % figdata: union_table

\textbf{4.} $e^{0.574} - 1 = 77.5~\%$.

\textbf{5.} It omits the \omterm{def:b3:gravity:multilateral-resistance}{multilateral resistance} terms. Large economies have many
nearby partners and low resistance, which lowers their bilateral trade with any
one partner relative to their size: the omitted term is correlated with GDP, and
its effect loads onto the GDP coefficients.

\textbf{6.} $-0.90$ against a median of $-0.89$ in the published estimates:
remarkably close. % ledger: b3gr:hm.t4

\textbf{7.} An exporter dummy absorbs everything specific to the exporter --- its
GDP and its resistance term --- and an importer dummy the same for the importer.

\textbf{8.} 19 for exporters and 19 for importers: 38 with a constant.

\textbf{9.} $(-0.979 + 1)/0.063 = 0.34$ standard errors: yes.

\textbf{10.} $0.750 \pm 1.96 \times 0.174$: from 0.41 to 1.09. It contains the true 0.5.

\textbf{11.} With a large partner close by, Brenland's resistance is low: its
alternatives are good. Relative to that, N10 is far away, and Brenland trades
less with it than with a remote country's equally distant partner.

\textbf{12.} A remote country has high resistance: all its alternatives are
costly, so a given bilateral cost is low relative to them, and two remote
countries trade more with each other than distance alone predicts.

\textbf{13.} \omterm{def:b3:gravity:ppml}{PPML} works with flows in levels, and a zero is a valid value of a
non-negative variable whose conditional mean is modelled.

\textbf{14.} $0.382 \pm 1.96 \times 0.170$: from 0.05 to 0.71. Partial effect $e^{0.382} - 1 =
46.5~\%$.

\textbf{15.} When the variance of the error falls with the size of the flow,
$E[\ln \eta \mid z]$ is more negative for small flows; small flows are the distant
and non-member pairs, so the log regression confounds the noise with distance
and with membership.

\textbf{16.} Distance: OLS with fixed effects ($-0.979$), then \omterm{def:b3:gravity:ppml}{PPML} ($-0.928$); the
union: \omterm{def:b3:gravity:ppml}{PPML} (0.382, against 0.750 with fixed effects by OLS and 0.574 by OLS);
both \omterm{def:b3:gravity:ppml}{PPML} errors are well within their standard errors.

\textbf{17.} $1 - e^{-0.382/5} = 7.3~\%$.

\textbf{18.} $e^{1.5/5} - 1 = 35.0~\%$.

\textbf{19.} The union lowers the members' \omterm{def:b3:gravity:multilateral-resistance}{multilateral resistance}: with cheaper
access to each other, each member's average trade cost falls, which offsets part
of the fall in their bilateral cost. Trade with members rises by 27.0~\%, not 46.5~\%.
% figdata: union_brenland

\textbf{20.} It falls by 10.5~\%: trade diversion, purchases switched from
non-members to members because members' goods became relatively cheaper.

\textbf{21.} Members' goods compete with Brenland's own at home; with Brenland's
output fixed and more of it exported, home sales fall, by 20.2~\%.

\textbf{22.} From $89.8 + 62.9 = 152.7$ to $114.1 + 56.3 = 170.4$: $+11.6~\%$.

\textbf{23.} No: gravity measures flows, not welfare. One needs the prices paid:
whether Brenland now buys from members goods that non-members supplied more
cheaply before (and loses the tariff revenue it collected on them), against the
gain from cheaper goods from members --- a model of welfare such as the gains
formula of \cref{ch:b3:ricardian-trade}.

\textbf{24.} In percent of trade with members, slightly larger: distant members
would be a small share of Brenland's trade, so its resistance would fall less,
and the rise would be closer to the 46.5~\% partial effect. But the union would act
on small flows, so the rise in Brenland's total trade, and the diversion away from
non-members, would be smaller.

\textbf{25.} The union raises Brenland's trade with its members by 27.0~\% and lowers
its trade with non-members by 10.5~\%.
\end{solution}
